\section{Independent Component Analysis}

\subsection{Motivation: Cocktail Party Problem}

% Cocktail Party Problem
\begin{frame}{Cocktail Party Problem}{}
	Zur Motivation betrachten wir das folgende Problem:
	
	\begin{itemize}
		\item Auf einer Party sprechen $N$ Personen gleichzeitig.
		\item Im Raum werden Mikrofone aufgestellt, welche \textit{überlappende Kombinationen} der Stimmen der Sprecher aufzeichnen.
		\item Der Einfachheit halber nehmen wir an, dass sich ebenfalls $N$ Mikrofone im Raum befinden.
		\item Da die Entfernung der Sprecher zu den einzelnen Mikrofonen verschieden ist, handelt es sich hierbei um leicht
			\textit{verschiedene Kombinationen}.
	\end{itemize}

	\textbf{Können wir die Stimmen der $N$ Sprecher aus den $N$ Aufzeichnungen isolieren?}
\end{frame}


% Beispiel: Cocktail Party Problem
\begin{frame}{Beispiel: Cocktail Party Problem}{}
	tbd
\end{frame}


\subsection{Formalisierung des Problems}

% Mathematische Formulierung des Problems
\begin{frame}{Mathematische Formulierung des Problems}{}
	\begin{itemize}
		\item Es seien die Signale $\bs(t) \in \R^N$ gegeben, welche von $N$ \textit{unabhängigen} (!) Quellen zum Zeitpunkt $t = 1, 2, \ldots, T$
			erzeugt werden (die Sprecher im Cocktail Party Problem).
		\item Allerdings beobachten wir nur die gemischten Daten $\bx(t)$ zum Zeitpunkt $t$ (die Aufnahmen der Mikrofone).
		\item Das ICA-Modell nimmt an, dass es sich hierbei um \textit{Linearkombinationen} der Quellsignale handelt.
		\item Wir betrachten also das Modell
		\begin{equation}
			\bx(t) = \bA \bs(t),
		\end{equation}
		wobei $\bA$ die \keyword{Mixing Matrix} ist.
	\end{itemize}
\end{frame}


% Mathematische Formulierung des Problems (Forts.)
\begin{frame}{Mathematische Formulierung des Problems (Forts.)}{}
	\begin{itemize}
		\item Unser Ziel ist es, die Quellsignale $\bs$ aus den beobachteten Mischsignalen $\bx$ zu rekonstruieren.
		\item Wäre $\bA$ bekannt, so könnten wir in den meisten Fällen $$\bs = \bA^{-1} \bx$$ berechnen und wären fertig.
		\item In der Regel ist $\bA$ unbekannt, wodurch das Problem deutlich schwerer wird.
		\item Wir definieren die \keyword{Unmixing Matrix} $\bB \coloneqq \bA^{-1}$.
		\item Mithilfe des ICA-Verfahrens werden wir diese Matrix approximieren. Mit ihr können wir dann die Quellsignale bestimmen, indem wir
			$\bs = \bB \bx$ berechnen.
	\end{itemize}
\end{frame}


\subsection{Annahmen und Limitationen von ICA}

% Annahmen des Verfahrens
\begin{frame}{Annahmen des Verfahrens}{}
	Die folgenden Annahmen sind wesentlich für das ICA-Verfahren:

	\begin{enumerate}
		\item \textbf{Zentrale Annahme:} Die Quellsignale (ICs: Independent Components) sind \textit{stochastisch unabhängig}.
		\item Die ICs sind \textit{nicht normalverteilt} (dazu später mehr).
		\item Die Mixing Matrix $\bA$ ist quadratisch. Es gibt allerdings Verallgemeinerungen der ICA, die auch rechteckige Matrizen erlauben.
			Wir machen diese Annahme hier, da das Problem selbst in diesem vereinfachten Fall schon schwer genug ist.
	\end{enumerate}
\end{frame}


% Limitationen des Verfahrens
\begin{frame}{Limitationen des Verfahrens}{}
	Es gibt eine Reihe von Einschränkungen bei der Rekonstruktion der Quellsignale:

	\begin{enumerate}
		\item Die ursprüngliche \textit{Reihenfolge} der Signale kann nicht rekonstruiert werden.
			Sei dazu $\bP$ eine \textit{Permutationsmatrix}. Dann ist
			$\bx = \bA \bs = \bA \bP^{-1} \bP \bs$, wobei $\bP \bs$ die permutierten Quellsignale sind.
			$\bA \bP^{-1}$ ist lediglich eine neue zu schätzende Mixing Matrix. Wir können folglich kein Signal \textit{das erste} Quellsignal nennen.
		\item Die \textit{Skalierung} der Quellsignale kann nicht wiederhergestellt werden
			(im Cocktail Party Problem entspricht dies der Lautstärke/Amplitude der Signale), wie man sich leicht überlegt: \\
			Sei dazu $\alpha \in \R$, dann gilt $$\bx = \bA \bs = \alpha \bA \cdot \frac{1}{\alpha} \bs.$$
		\item Normalverteilte Signale können nicht aufgetrennt werden.
	\end{enumerate}
\end{frame}


% Warum sind Normalverteilungen verboten?
\begin{frame}{Warum sind Normalverteilungen verboten?}{}
	\begin{itemize}
		\item Wir betrachten ein Beispiel mit $N = 2$ und $\bs \sim N_{\0, \bI}$.
		\item Die Konturlinien der Standardnormalverteilung sind perfekte Kreise, die auf dem Koordinatenursprung zentriert sind.
		\item Insbesondere ist die Dichte \keyword{rotationsinvariant}.
		\item Für $\bx = \bA \bs$ gilt damit: $\bx \sim N_{\0, \bA \bA^\top}$. \\ \vspace*{2mm}
		\begin{myproof}
			Mit der Linearität des Erwartungswerts gilt $\E_{\bs}[\bx] = \E_{\bs}[\bA \bs] = \bA \E_{\bs}[s] = \0$.
			Die Streuungsmatrix ist
			\begin{align*}
				\D[\bx]
					&= \E_{\bs}\bigl[ \bx \bx^\top \bigr]
						= \E_{\bs}\bigl[ \bA \bs \bs^\top \bA^\top \bigr]
						= \bA \E\bigl[ \bs \bs^\top \bigr] \bA^\top
						= \bA \cdot \D[\bs] \cdot \bA^\top
			\\
					&= \bA \bA^\top.
			\end{align*}
			Also gilt die Behauptung.
		\end{myproof} 
	\end{itemize}
\end{frame}


% Warum sind Normalverteilungen verboten? (Forts.)
\begin{frame}{Warum sind Normalverteilungen verboten? (Forts.)}{}
	\begin{itemize}
		\item Sei nun $\bR \in \R^{N \times N}$ eine beliebige \textit{orthogonale Matrix}. \\ \vspace*{2mm}
		\textbf{Eigenschaften orthogonaler Matrizen:}
		\begin{itemize}
			\item Orthogonale Matrizen beschreiben eine Drehung/Spiegelung im Raum.
			\item Es gilt stets $\bR \bR^\top = \bR^\top \bR = \bI.$
			\item Es ist $\det \bR = 1$ für alle orthogonalen Matrizen $\bR$.
			\item Da $\bR$ nur eine Drehung beschreibt, gilt für alle $\bx \in \R^N$ die Identität $\Vert \bR \bx \Vert = \Vert \bx \Vert$.
		\end{itemize}
		\item Sei ferner $\bm{\tilde{A}} \coloneqq \bA \bR$ die Matrix, welche mit der von $\bR$ beschriebenen Drehung aus $\bA$ hervorgeht.
		\item Dann hat der Vektor $\bm{\tilde{x}} = \bm{\tilde{A}} \bs$ dieselbe Verteilung wie $\bx$. \\ \vspace*{2mm}
		\begin{myproof}
			Den Erwartungswert rechnet man analog zu oben aus. Aufgrund der Orthogonalität von $\bR$ gilt
			\begin{equation*}
				\E_{\bs}\bigl[ \bm{\tilde{x}} \bm{\tilde{x}}^\top \bigr]
					= \E_{\bs}\bigl[ \bm{\tilde{A}} \bs \bs^\top \bm{\tilde{A}}^\top \bigr]
					= \E_{\bs}\bigl[ \bA \bR \bs \bs^\top \bR^\top \bA^\top \bigr]
					= \bA \bR \cdot \D[\bs] \cdot \bR^\top \bA^\top
					= \bA \bR \bR^\top \bA
					= \bA \bA^\top,
			\end{equation*}
			sodass $\bm{\tilde{x}} \sim N_{\0, \bA \bA^\top}$ wie behauptet gilt.
		\end{myproof}
	\end{itemize}
\end{frame}


% Warum sind Normalverteilungen verboten? (Forts.)
\begin{frame}{Warum sind Normalverteilungen verboten? (Forts.)}{}
	\begin{itemize}
		\item Die voranstehenden Ausführungen zeigen, können wir in diesem Fall nicht rekonstruieren, ob die Daten mit der Matrix
			$\bA$ oder $\bm{\tilde{A}}$ gemixt wurden.
		\item Für normalverteilte Daten können wir folglich die Komponenten \textit{nur bis auf eine Drehung rekonstruieren}.
	\end{itemize}
\end{frame}


\subsection{Visualisierung der ICA}

% Ein einfaches Beispiel
\begin{frame}{Ein einfaches Beispiel}{}

\end{frame}

